\documentclass[12pt]{article}
\usepackage{amsmath,amssymb}

\begin{document}
\section*{{2024 AMC 12A Problems/Problem 8}}
Source: AoPS Wiki

\subsection*{Solution}
Note that this is equivalent to $\sin(3\theta)\cos(2\theta)=1$ , which is clearly only possible when $\sin(3\theta)=\cos(2\theta)=\pm1$ . (If either one is between $1$ and $-1$ , the other one must be greater than $1$ or less than $-1$ to offset the product, which is impossible for sine and cosine.) They cannot be both $-1$ since we cannot take logarithms of negative numbers, so they are both $+1$ . Then $3\theta$ is $\dfrac\pi2$ more than a multiple of $2\pi$ and $2\theta$ is a multiple of $2\pi$ , so $\theta$ is $\dfrac\pi6$ more than a multiple of $\dfrac23\pi$ and also a multiple of $\pi$ . However, a multiple of $\dfrac23\pi$ will always have a denominator of $1$ or $3$ , and never $6$ ; it can thus never add with $\dfrac\pi6$ to form an integral multiple of $\pi$ . Thus, there are $\boxed{\textbf{(A) }0}$ solutions. ~Technodoggo

\end{document}
